% Drawing 161-6/02-1 -- folder 161-6, batch 02, archive page 27 (PDF page 28).
% Transcription: croquis en marge gauche, rattaché à « NB » par un trait : deux flèches parallèles $G_\rho \rightrightarrows G_\rho$ marquées $\bar f, \bar g$, une flèche montante de $\pi_1(X)$ vers le premier $G_\rho$, et une courbe sans tête de $\pi_1(X)$ au second
% Kind: margin arrow diagram. Method 2 (schematic redraw from the facsimile). Status: ok.
% Notes: Two parallel arrows G_rho => G_rho marked bar f, bar g, an arrow going up from pi_1(X) to the first G_rho, and a headless curve from pi_1(X) to the second. The stroke tying the sketch to « NB » passes through the crop (vertical stroke above the second G_rho, horizontal stroke to its right); only its part inside the crop is drawn, in pencil. The bars of bar f and bar g are written as two short strokes above the letters.
% Coordinates: 1 unit = 1 mm = 1/100 of the crop width (bbox in page fractions: [0.14, 0.65, 0.26, 0.76]).
% Redrawn by Claude (claude-opus-5-5), issue #73, 2026-11-07; not checked by a human; not compiled (no TeX in the session).
% Rights: reproduces part of the work -- behind RENDER_TRANSCRIPTS, never published by itself.
\documentclass[tikz,border=3pt]{standalone}
\usepackage{amsmath,amssymb}
\usetikzlibrary{patterns,arrows.meta}
\tikzset{pencil/.style={black!45}, struck/.style={black!35}, mark/.style={circle,draw,inner sep=1.3pt}, pt/.style={circle,fill,inner sep=0.9pt}, opt/.style={circle,draw,fill=white,inner sep=1.1pt}}
\begin{document}
\begin{tikzpicture}[x=1mm,y=1mm,line cap=round,line join=round,font=\scriptsize]
\node[inner sep=1pt] at (26,-50) {\scriptsize $G_\rho$};
\node[inner sep=1pt] at (85,-52) {\scriptsize $G_\rho$};
\node[inner sep=1pt] at (50,-29) {\tiny $\bar f,\bar g$};
\draw[-{Stealth[length=2mm]}] (41,-47.5) -- (63,-47.5);
\draw[-{Stealth[length=2mm]}] (41,-52.5) -- (64,-52.5);
\draw[-{Stealth[length=2mm]},very thick] (24.5,-102) -- (23,-68);
\node[inner sep=1pt] at (27,-114) {\tiny $\pi_1(X)$};
\draw plot[smooth, tension=0.55] coordinates {(40,-113.5) (60,-110) (74,-101) (80,-87) (83,-68)};
\draw[pencil] (77,0) -- (77,-53);
\draw[pencil] plot[smooth, tension=0.55] coordinates {(80,-59) (88,-62.5) (100,-63.5)};
\end{tikzpicture}
\end{document}
